\documentclass[12pt,a4paper]{article}

% ── Packages ──────────────────────────────────────────────────────────────────
\usepackage[T1]{fontenc}
\usepackage[utf8]{inputenc}
\usepackage[english]{babel}
\usepackage{geometry}
\geometry{top=2.5cm, bottom=2.5cm, left=2.5cm, right=2.5cm}

\usepackage{graphicx}
\usepackage{xcolor}
\usepackage{amsmath}
\usepackage{amssymb}
\usepackage{booktabs}
\usepackage{array}
\usepackage{longtable}
\usepackage{enumitem}
\usepackage{hyperref}
\usepackage{fancyhdr}
\usepackage{titlesec}
\usepackage{tcolorbox}
\usepackage{parskip}

% ── Colours ───────────────────────────────────────────────────────────────────
\definecolor{primaryblue}{RGB}{0,123,255}
\definecolor{darkblue}{RGB}{0,70,140}
\definecolor{lightgray}{RGB}{248,249,250}
\definecolor{successgreen}{RGB}{21,87,36}
\definecolor{successbg}{RGB}{212,237,218}
\definecolor{warningyellow}{RGB}{133,100,4}
\definecolor{warningbg}{RGB}{255,243,205}
\definecolor{dangerbg}{RGB}{248,215,218}
\definecolor{dangerred}{RGB}{114,28,36}
\definecolor{codebg}{RGB}{233,236,239}

% ── Hyperref setup ────────────────────────────────────────────────────────────
\hypersetup{
    colorlinks=true,
    linkcolor=darkblue,
    urlcolor=primaryblue,
    citecolor=darkblue,
    pdftitle={Fatigue 3D Model – User Manual},
    pdfauthor={Statistical Analysis Tool},
    pdfsubject={User Manual}
}

% ── Section formatting ────────────────────────────────────────────────────────
\titleformat{\section}{\Large\bfseries\color{darkblue}}{}{0em}{}[\color{primaryblue}\titlerule]
\titleformat{\subsection}{\large\bfseries\color{darkblue}}{}{0em}{}
\titleformat{\subsubsection}{\normalsize\bfseries\color{darkblue}}{}{0em}{}

% ── Custom boxes ──────────────────────────────────────────────────────────────
\tcbuselibrary{skins,breakable}

\newtcolorbox{notebox}{
    colback=successbg, colframe=successgreen,
    fonttitle=\bfseries, title=Note,
    breakable, left=6pt, right=6pt, top=4pt, bottom=4pt
}

\newtcolorbox{warnbox}{
    colback=warningbg, colframe=warningyellow,
    fonttitle=\bfseries, title=Warning,
    breakable, left=6pt, right=6pt, top=4pt, bottom=4pt
}

\newtcolorbox{stepbox}[1]{
    colback=lightgray, colframe=primaryblue,
    fonttitle=\bfseries\color{white}, coltitle=white,
    title=Step #1,
    breakable, left=6pt, right=6pt, top=4pt, bottom=4pt
}

\newtcolorbox{infobox}{
    colback=codebg, colframe=primaryblue!60,
    breakable, left=6pt, right=6pt, top=4pt, bottom=4pt,
    fontupper=\ttfamily\small
}

% ── Header / footer ───────────────────────────────────────────────────────────
\pagestyle{fancy}
\fancyhf{}
\fancyhead[L]{\small\color{darkblue}\textbf{Fatigue 3D Model}}
\fancyhead[R]{\small\color{darkblue}User Manual}
\fancyfoot[C]{\small\thepage}
\renewcommand{\headrulewidth}{0.4pt}
\renewcommand{\headrule}{\color{primaryblue}\hrule width\headwidth height\headrulewidth}

% ── Inline code ───────────────────────────────────────────────────────────────
\newcommand{\code}[1]{\colorbox{codebg}{\texttt{\small#1}}}

% ══════════════════════════════════════════════════════════════════════════════
\begin{document}
% ══════════════════════════════════════════════════════════════════════════════

% ── Title page ────────────────────────────────────────────────────────────────
\begin{titlepage}
    \centering
    \vspace*{2cm}
    {\color{darkblue}\rule{\linewidth}{2pt}}\\[0.5cm]
    {\Huge\bfseries\color{darkblue} Fatigue 3D Model\\[0.4cm]
     \Large Statistical Analysis Tool}\\[0.3cm]
    {\color{primaryblue}\rule{\linewidth}{2pt}}\\[1.5cm]
    {\LARGE\bfseries User Manual}\\[0.5cm]
    {\large Version 1.0}\\[3cm]
    \begin{tcolorbox}[colback=lightgray, colframe=primaryblue, width=0.8\linewidth]
        \centering
        \textbf{Program:} \texttt{Fatigue3DModel.html}\\[4pt]
        \textbf{Purpose:} Three-dimensional statistical modelling\\
        \phantom{\textbf{Purpose:}} of fatigue data using hyperbolic regression\\[4pt]
        \textbf{Platform:} Any modern web browser\\[4pt]
        \textbf{Input format:} CSV (comma-separated values)
    \end{tcolorbox}
    \vfill
    {\small\color{gray} This document describes the complete workflow of the tool,\\
    from data loading to the generation of percentile hyperbola plots.}
\end{titlepage}

% ── Table of contents ─────────────────────────────────────────────────────────
\tableofcontents
\newpage

% ══════════════════════════════════════════════════════════════════════════════
\section{Introduction}
% ══════════════════════════════════════════════════════════════════════════════

\textbf{Fatigue 3D Model} is a self-contained, browser-based statistical analysis tool
designed to fit a three-dimensional hyperbolic regression model to fatigue or similar
engineering data.  The program is distributed as a single HTML file
(\texttt{Fatigue3DModel.html}) and requires no installation: it runs entirely inside
any modern web browser.

\subsection{Purpose}

The tool addresses the problem of characterising the statistical distribution of a
\emph{response variable} (e.g.\ fatigue life) as a function of two \emph{predictor
variables} (e.g.\ stress amplitude and a third physical quantity such as stress ratio,
temperature, or frequency).  By fitting a multivariate linear model in the transformed
space $x \cdot y \cdot z$ and computing empirical percentiles of the residuals, the
tool generates a family of hyperbolic iso-percentile curves that cover the full
probability range from $P_0$ to $P_{99}$.

\subsection{Key Features}

\begin{itemize}[leftmargin=1.5em]
    \item Automatic detection and application of a base-10 logarithmic transformation
          for columns whose range spans more than three orders of magnitude.
    \item Interactive scatter plot of the raw data with click-to-toggle outlier
          identification.
    \item Ordinary least-squares (OLS) fitting of a seven-parameter bilinear model
          in the product space.
    \item Automatic outlier detection based on the asymptotes of the fitted hyperbola.
    \item Computation of residual quantiles at seven probability levels:
          $P_0$, $P_5$, $P_{25}$, $P_{50}$, $P_{75}$, $P_{95}$, and $P_{99}$.
    \item Generation of iso-percentile hyperbola plots for five representative values
          of the third variable, both in original scale and in the uniform $[0,1]$
          probability scale.
    \item One-click PDF export of every interactive plot.
\end{itemize}

\subsection{System Requirements}

\begin{itemize}[leftmargin=1.5em]
    \item A modern web browser (Chrome, Firefox, Edge, or Safari --- version released
          after 2020 recommended).
    \item An internet connection is required on first use so that the browser can load
          the external libraries (Plotly, jsPDF, html2canvas) from their CDN addresses.
          Subsequent uses may work offline if the browser has cached those resources.
    \item No installation, no server, no database.
\end{itemize}

% ══════════════════════════════════════════════════════════════════════════════
\section{Input Data Format}
% ══════════════════════════════════════════════════════════════════════════════

The program reads a \textbf{comma-separated values (CSV)} file.  The file must satisfy
the following requirements.

\subsection{File Structure}

\begin{itemize}[leftmargin=1.5em]
    \item \textbf{First row:} column headers, separated by commas.
    \item \textbf{Subsequent rows:} numeric data, one observation per row.
    \item \textbf{Missing values:} leave the cell empty or enter a non-numeric token;
          the program will treat it as \texttt{null} and skip that observation for
          any computation involving the missing column.
    \item \textbf{Encoding:} UTF-8 plain text.
    \item \textbf{Decimal separator:} period (\texttt{.}).
\end{itemize}

\subsection{Required Columns}

The tool expects \textbf{exactly three numeric columns}.  Their roles (horizontal
variable $X$, vertical variable $Y$, and third variable $Z$) are assigned interactively
after loading, so the physical order in the file is arbitrary.

\begin{center}
\begin{tabular}{@{}llp{7cm}@{}}
\toprule
\textbf{Column} & \textbf{Role} & \textbf{Typical example (fatigue)} \\
\midrule
1\textsuperscript{st} & $X$ --- horizontal predictor & Stress amplitude $\sigma_a$ \\
2\textsuperscript{nd} & $Y$ --- vertical response     & Fatigue life $N$ \\
3\textsuperscript{rd} & $Z$ --- third predictor        & Stress ratio $R$ \\
\bottomrule
\end{tabular}
\end{center}

\subsection{Automatic Logarithmic Transformation}

At load time the program inspects each column.  If the ratio
$\max / \min > 1000$ and all values are strictly positive, the entire column is
replaced by $\log_{10}(\text{value})$ before any further computation.  The column
header is then labelled \texttt{log(OriginalName)} throughout the interface.

\begin{notebox}
    The transformation is applied silently.  After loading, check the information
    panel to confirm which columns have been log-transformed.
\end{notebox}

\subsection{Minimal Example}

\begin{infobox}
sigma\_a,N,R\\
200,1.2e6,0.1\\
250,4.5e5,0.1\\
300,1.8e5,0.1\\
200,9.8e5,-0.5\\
250,3.7e5,-0.5
\end{infobox}

% ══════════════════════════════════════════════════════════════════════════════
\section{Program Interface Overview}
% ══════════════════════════════════════════════════════════════════════════════

The interface is divided into six numbered sections that follow the natural analysis
workflow.  They are presented sequentially from top to bottom on the page.

\begin{center}
\renewcommand{\arraystretch}{1.3}
\begin{tabular}{@{}clp{8.5cm}@{}}
\toprule
\textbf{Part} & \textbf{Title} & \textbf{Purpose} \\
\midrule
I   & Data Reading and Variable Definition & Load the CSV file and choose which column plays each role. \\
II  & Data Table & Inspect the data in original or $[0,1]$ scale; identify outliers visually. \\
III & Three-dimensional Model Fitting & Fit the OLS model and inspect residuals. \\
IV  & (internal)  & Outlier detection (executed as part of Part III). \\
V   & Percentile Calculation & Display the seven residual quantiles. \\
VI  & Hyperbolas for Different $Z$ Values & Generate and export the iso-percentile curves. \\
\bottomrule
\end{tabular}
\end{center}

% ══════════════════════════════════════════════════════════════════════════════
\section{Step-by-Step Workflow}
% ══════════════════════════════════════════════════════════════════════════════

\subsection{Part I --- Data Loading and Variable Definition}

\begin{stepbox}{1 --- Open the tool}
Open \texttt{Fatigue3DModel.html} in your web browser by double-clicking the file
(local mode) or by navigating to the hosted URL (GitHub Pages mode).  The page
detects its environment automatically.
\end{stepbox}

\begin{stepbox}{2 --- Set the bandwidth parameter}
The \textbf{var3bandwidth} input (default \texttt{0.10}) controls the neighbourhood
size used when computing local statistics for a given $Z$ value.  A value of
\texttt{0.10} means that data points within $\pm 10\%$ of the full $Z$-range around
the target value are included.  Acceptable values are strictly between 0 and 1.

Click \textbf{Apply var3bandwidth} to confirm the value.  A green confirmation
message will appear.
\end{stepbox}

\begin{stepbox}{3 --- Load the CSV file (local mode)}
Click \textbf{Choose File}, select your \texttt{.csv} file, then click
\textbf{Load File}.  A success message confirms the number of rows and columns
detected, together with a list of any log-transformed columns.
\end{stepbox}

\begin{notebox}
    When the tool is served from GitHub Pages it automatically loads a file named
    \texttt{DataFatigue.csv} located in the repository root.  In that case the
    file-picker control is hidden.
\end{notebox}

\begin{stepbox}{4 --- Assign variables}
Two drop-down menus appear after a successful load:
\begin{itemize}[leftmargin=1.5em]
    \item \textbf{Vertical Variable (Y):} the response to be plotted on the vertical axis.
    \item \textbf{Horizontal Variable (X):} the predictor on the horizontal axis.
\end{itemize}
The remaining column is automatically designated as the \textbf{third variable} ($Z$).
Click \textbf{Apply Selection}.  An interactive scatter plot of $Y$ vs.\ $X$ is
displayed immediately below.
\end{stepbox}

\begin{warnbox}
    You must select two \emph{different} columns for $X$ and $Y$.  The tool will
    alert you if the same column is chosen for both.
\end{warnbox}

\subsection{Part II --- Data Table}

The data table shows all observations.  Rows highlighted in red correspond to
observations currently classified as outliers.

The toggle button \textbf{Show Data in [0,1] Scale} / \textbf{Show Original Data}
switches between the original (possibly log-transformed) values and the empirical
CDF-based probability scale.  The $[0,1]$ scale is computed column-by-column using
linear interpolation of the empirical CDF:
\[
    u_i = \frac{\text{rank}(x_i)}{n}, \qquad i = 1, \dots, n.
\]

\subsection{Part III --- Three-dimensional Model Fitting}

Click \textbf{Fit Model and Analyze}.  The program executes the following sequence
automatically:

\subsubsection{Model Specification}

The tool fits a seven-parameter ordinary least-squares model.  Define
$x = X_i$, $y = Y_i$, $z = Z_i$ for the $i$-th observation.  The model is
linearised by expanding the product $xyz$ and expressing the response as:

\[
    -\,x_i y_i z_i \;=\;
    C_0
    + C_1 x_i
    + C_2 y_i
    + C_3 z_i
    + C_{12}\,x_i y_i
    + C_{13}\,x_i z_i
    + C_{23}\,y_i z_i
    + \varepsilon_i
\]

This can be written compactly as $\mathbf{y} = \mathbf{X}\boldsymbol{\theta} + \boldsymbol{\varepsilon}$,
where the design matrix $\mathbf{X}$ has columns
$[1,\, x,\, y,\, z,\, xy,\, xz,\, yz]$ and the parameter vector is
$\boldsymbol{\theta} = [C_0, C_1, C_2, C_3, C_{12}, C_{13}, C_{23}]^\top$.

The OLS estimator $\hat{\boldsymbol{\theta}} = (\mathbf{X}^\top\mathbf{X})^{-1}\mathbf{X}^\top\mathbf{y}$
is solved using Gaussian elimination with partial pivoting.

\subsubsection{Goodness-of-Fit Statistics}

After fitting, the following statistics are displayed:

\begin{center}
\renewcommand{\arraystretch}{1.3}
\begin{tabular}{@{}ll@{}}
\toprule
\textbf{Statistic} & \textbf{Definition} \\
\midrule
Mean of residuals  & $\bar{\varepsilon} = \frac{1}{n}\sum_i \varepsilon_i$ \\
Std.\ deviation    & $s = \sqrt{\frac{1}{n-1}\sum_i (\varepsilon_i - \bar{\varepsilon})^2}$ \\
Variance           & $s^2$ \\
MSE                & $\frac{1}{n}\sum_i \varepsilon_i^2$ \\
RMSE               & $\sqrt{\text{MSE}}$ \\
MAE                & $\frac{1}{n}\sum_i |\varepsilon_i|$ \\
$R^2$              & $1 - \text{SS}_\text{res}/\text{SS}_\text{tot}$ \\
MAPE               & $\frac{100}{n}\sum_i |\varepsilon_i / y_i|$ (when $y_i \neq 0$) \\
\bottomrule
\end{tabular}
\end{center}

\subsubsection{Residual Empirical CDF Plot}

A plot of the empirical cumulative distribution function of the residuals
$F(\varepsilon)$ is generated.  The seven percentile values are overlaid as
coloured markers.  The plot can be exported to PDF via the
\textbf{Download Plot as PDF} button.

\subsubsection{Automatic Outlier Detection}

Outlier detection is carried out in one pass over five equally spaced $Z$ reference
values spanning the observed range $[Z_\min, Z_\max]$.  For each reference value
$z_k$ the fitted model defines a rectangular hyperbola with:
\begin{itemize}[leftmargin=1.5em]
    \item a \textbf{vertical asymptote} at $X = X_\text{asym}(z_k)$,
    \item a \textbf{horizontal asymptote} at $Y = Y_\text{asym}(z_k)$.
\end{itemize}
Any observation within the bandwidth of $z_k$ that lies to the \emph{left} of the
vertical asymptote or \emph{below} the horizontal asymptote is flagged as an outlier
and removed from subsequent computations.  The scatter plot is refreshed and outlier
rows are highlighted in the data table.

\begin{warnbox}
    At least \textbf{7 valid observations} must remain after outlier removal for the
    analysis to continue.  The program will halt and display a warning if this
    condition is not met.
\end{warnbox}

\subsection{Part V --- Percentile Calculation}

Seven percentiles of the global residual distribution are computed and displayed:

\begin{center}
$P_0,\quad P_5,\quad P_{25},\quad P_{50},\quad P_{75},\quad P_{95},\quad P_{99}$
\end{center}

These are obtained by linear interpolation between the order statistics of the
sorted residual vector.  The quantile at probability level $p$ is:

\[
    Q(p) = (1-w)\,\varepsilon_{(j)} + w\,\varepsilon_{(j+1)},
    \qquad j = \lfloor p(n-1) \rfloor,\quad w = p(n-1) - j.
\]

\subsection{Part VI --- Hyperbola Plots}

\subsubsection{Hyperbola Generation}

For each of five reference values $z_k$ (equally spaced in $[Z_\min, Z_\max]$),
the program computes seven iso-percentile hyperbolic curves.  Given a quantile
$q$ and fixed $z_k$, the $Y$ value corresponding to a given $X$ is obtained
by solving the model equation for $Y$:

\[
    Y(X;\, z_k,\, q) = \frac{q - C_0 - C_1 X - C_3 z_k - C_{13} X z_k}
                            {C_2 + C_{12} X + C_{23} z_k}
\]

where $q$ replaces the residual $\varepsilon$ term.  Only points that lie in the
physically valid region (beyond both asymptotes) are plotted.

\subsubsection{Original-Scale Plots}

For each reference $Z$ value a two-panel row is displayed:
\begin{itemize}[leftmargin=1.5em]
    \item \textbf{Left panel:} the seven iso-percentile hyperbolas overlaid on the
          scatter plot.  Data points whose $Z$ value falls within the bandwidth
          are shown in orange; outliers in blue; remaining data in light grey.
    \item \textbf{Right panel:} the empirical CDF of the local residuals (those
          belonging to the bandwidth neighbourhood of $z_k$), with the seven
          percentile markers.
\end{itemize}

A line of text below each row lists the local percentile values for that $Z$ slice.

\subsubsection{Plots in \texorpdfstring{$[0,1]$}{[0,1]} Probability Scale}

Click \textbf{Generate Plots in [0,1] Scale} to produce an additional set of plots
where both axes are transformed to the uniform probability scale using the empirical
CDF:

\[
    u_X = F_X(X), \qquad u_Y = F_Y(Y).
\]

In this representation the axes run from 0 to 1 and all variables are treated
symmetrically, facilitating visual comparison across different datasets.

\subsubsection{Colour Convention}

All hyperbola plots use the following consistent colour scheme:

\begin{center}
\renewcommand{\arraystretch}{1.2}
\begin{tabular}{@{}clc@{}}
\toprule
\textbf{Percentile} & \textbf{Colour} & \textbf{Line style} \\
\midrule
$P_0$   & Red (\texttt{\#FF0000})      & Solid \\
$P_5$   & Orange-red (\texttt{\#FF4500}) & Solid \\
$P_{25}$ & Dark orange (\texttt{\#FF8C00}) & Solid \\
$P_{50}$ & Orange (\texttt{\#FFA500})    & \textbf{Dashed} \\
$P_{75}$ & Steel blue (\texttt{\#4682B4}) & Solid \\
$P_{95}$ & Dodger blue (\texttt{\#1E90FF}) & Solid \\
$P_{99}$ & Royal blue (\texttt{\#4169E1}) & Solid \\
\bottomrule
\end{tabular}
\end{center}

% ══════════════════════════════════════════════════════════════════════════════
\section{Exporting Results}
% ══════════════════════════════════════════════════════════════════════════════

Every plot generated by the tool can be exported individually as a PDF file.
Clicking any \textbf{Download Plot as PDF} button (shown in red) performs the
following sequence internally:

\begin{enumerate}[leftmargin=1.5em]
    \item The plot is temporarily resized to $1200 \times 900$ pixels with a white
          background.
    \item A PNG image is captured at $2\times$ pixel density ($2400 \times 1800$ px).
    \item The image is embedded in a landscape A4 PDF page using the jsPDF library.
    \item The PDF is downloaded through the browser under an automatically generated
          filename (e.g.\ \texttt{Hyperbola\_0.250.pdf} or
          \texttt{Hyperbola\_01\_0.250.pdf} for the $[0,1]$-scale version).
    \item The plot is restored to its original display size.
\end{enumerate}

% ══════════════════════════════════════════════════════════════════════════════
\section{Mathematical Background}
% ══════════════════════════════════════════════════════════════════════════════

\subsection{The Three-Variable Hyperbolic Model}

The fundamental relationship exploited by the tool is that, for fixed $z$, the
implicit equation

\[
    C_0 + C_1 x + C_2 y + C_3 z + C_{12}\,xy + C_{13}\,xz + C_{23}\,yz
    \;=\; -\,xyz
\]

defines a rectangular hyperbola in the $(x, y)$ plane.  Rearranging:

\[
    y\,(C_2 + C_{12}\,x + C_{23}\,z)
    \;=\; -xyz - C_0 - C_1 x - C_3 z - C_{13}\,xz
\]
\[
    y \;=\; \frac{-xyz - C_0 - C_1 x - C_3 z - C_{13}\,xz}
                 {C_2 + C_{12}\,x + C_{23}\,z}.
\]

The \textbf{vertical asymptote} is the value of $x$ for which the denominator
(or an equivalent grouping) becomes zero, and the \textbf{horizontal asymptote}
follows symmetrically.

\subsection{Iso-Percentile Curves}

Adding a residual quantile $q$ to the right-hand side shifts the hyperbola
vertically in the product space:

\[
    C_0 + C_1 x + C_2 y + C_3 z + C_{12}\,xy + C_{13}\,xz + C_{23}\,yz + q
    \;=\; -\,xyz.
\]

Each value of $q$ (i.e.\ each of the seven percentile levels) produces a distinct
hyperbola.  Together they form the \emph{iso-percentile fan} displayed in Part VI.

\subsection{Empirical CDF and \texorpdfstring{$[0,1]$}{[0,1]} Scale}

The empirical cumulative distribution function is estimated by the \textbf{Weibull
plotting position}:

\[
    \hat{F}(x_{(i)}) = \frac{i}{n}, \qquad i = 1, \dots, n,
\]

where $x_{(1)} \le x_{(2)} \le \cdots \le x_{(n)}$ are the order statistics.
Linear interpolation is used between adjacent pairs to map any value onto
$[0,1]$ and, conversely, to invert from $[0,1]$ back to the original scale.

% ══════════════════════════════════════════════════════════════════════════════
\section{Troubleshooting}
% ══════════════════════════════════════════════════════════════════════════════

\renewcommand{\arraystretch}{1.4}
\begin{longtable}{@{}p{5cm}p{9cm}@{}}
\toprule
\textbf{Symptom} & \textbf{Likely cause and remedy} \\
\midrule
\endhead
\bottomrule
\endfoot
\textit{``Plotly has not loaded''} red banner on startup &
    The browser could not reach the Plotly CDN.  Check your internet connection and
    reload the page. \\
\textit{``Please select a CSV file''} alert when clicking Load &
    No file was selected in the file picker.  Click the picker first. \\
File loads but columns are wrong / missing &
    Verify that the CSV uses commas as delimiters, that the first row contains
    column names, and that all data values use a period as the decimal separator. \\
\textit{``Not enough valid data for fitting (minimum 7 points)''} &
    The dataset has fewer than 7 rows with complete (non-null, finite) values in
    all three columns.  Check for missing values or remove rows with \texttt{NaN}. \\
\textit{``Singular matrix''} error during fitting &
    Two or more columns (after log-transformation) are linearly dependent.  This
    can happen if one variable is a constant or a perfect linear function of
    another. \\
Hyperbola curves appear truncated or absent &
    The curves may fall entirely outside the visible axis range.  Try a different
    variable assignment or verify that the bandwidth parameter captures enough
    local data points. \\
PDF download button is greyed out &
    The corresponding plot has not been generated yet.  Complete the preceding
    analysis step first. \\
\textit{``No valid data for plot''} message in scatter area &
    All rows contain at least one null or non-finite value in the selected
    variables.  Inspect the raw CSV for formatting issues. \\
\end{longtable}

% ══════════════════════════════════════════════════════════════════════════════
\section{Parameter Reference}
% ══════════════════════════════════════════════════════════════════════════════

\begin{center}
\renewcommand{\arraystretch}{1.4}
\begin{tabular}{@{}lp{4cm}p{7cm}@{}}
\toprule
\textbf{Parameter} & \textbf{Location} & \textbf{Description} \\
\midrule
\texttt{var3bandwidth} &
    Part I control panel &
    Fractional half-width of the $Z$-neighbourhood used for local statistics.
    Range: $(0,1)$.  Default: \texttt{0.10}. \\
$C_0,\dots,C_{23}$ &
    Displayed in Part III after fitting &
    Seven OLS regression coefficients of the linearised hyperbolic model. \\
$P_0, P_5, P_{25}, P_{50}, P_{75}, P_{95}, P_{99}$ &
    Displayed in Parts III, V, VI &
    Empirical quantiles of the global residual distribution. \\
\texttt{ecdfX}, \texttt{ecdfY} &
    Internal (used in Part VI) &
    Empirical CDF lookup tables for the $X$ and $Y$ columns, used for the
    $[0,1]$-scale transformation. \\
\bottomrule
\end{tabular}
\end{center}

% ══════════════════════════════════════════════════════════════════════════════
\section{Quick-Start Checklist}
% ══════════════════════════════════════════════════════════════════════════════

\begin{enumerate}[leftmargin=1.5em, label=\textbf{\arabic*.}]
    \item Open \texttt{Fatigue3DModel.html} in your browser.
    \item Confirm the green \textit{``Plotly loaded successfully''} banner.
    \item (Optional) Adjust \texttt{var3bandwidth} and click \textbf{Apply var3bandwidth}.
    \item Select your CSV file and click \textbf{Load File}.
    \item Verify the column list and log-transformation notice.
    \item Select the \textbf{Vertical} and \textbf{Horizontal} variables; click
          \textbf{Apply Selection}.
    \item Inspect the scatter plot; use \textbf{Show Data in [0,1] Scale} if desired.
    \item Click \textbf{Fit Model and Analyze}.
    \item Read the parameter estimates and goodness-of-fit metrics in Part III.
    \item Acknowledge the outlier-detection alert (lists removed points).
    \item Inspect the residual CDF plot and the $P$-values in Part V.
    \item Scroll to Part VI to view the iso-percentile hyperbola plots.
    \item Click \textbf{Generate Plots in [0,1] Scale} for the normalised view.
    \item Use the red \textbf{Download Plot as PDF} buttons to export any plot.
\end{enumerate}

% ══════════════════════════════════════════════════════════════════════════════
\section{Glossary}
% ══════════════════════════════════════════════════════════════════════════════

\begin{description}[leftmargin=2em, style=nextline]
    \item[Bandwidth (\texttt{var3bandwidth})]
        The half-width of the sliding window in the $Z$-direction, expressed as
        a fraction of the total $Z$-range.
    \item[ECDF]
        Empirical Cumulative Distribution Function.  A step function that estimates
        the true CDF from the sample.
    \item[Horizontal asymptote]
        The value $Y_\text{asym}$ towards which the fitted hyperbola converges as
        $X \to \infty$.
    \item[Iso-percentile curve]
        A hyperbola in the $(X, Y)$ plane along which the probability that a random
        observation falls below the curve is constant.
    \item[Log-transformation]
        Replacement of raw values by their base-10 logarithm, applied automatically
        to columns with a dynamic range exceeding three orders of magnitude.
    \item[MAE]
        Mean Absolute Error: $\frac{1}{n}\sum |\varepsilon_i|$.
    \item[MAPE]
        Mean Absolute Percentage Error: $\frac{100}{n}\sum |\varepsilon_i / y_i|$.
    \item[MSE / RMSE]
        Mean Squared Error / Root Mean Squared Error.
    \item[OLS]
        Ordinary Least Squares regression.
    \item[Outlier]
        A data point that lies on the wrong side of a fitted hyperbola's asymptote,
        inconsistent with the physical model.
    \item[Percentile $P_k$]
        The value below which $k\%$ of the empirical residual distribution falls.
    \item[Residual $\varepsilon_i$]
        Difference between the model prediction and the observed response for
        observation $i$ in the transformed product space.
    \item[Third variable ($Z$)]
        The conditioning variable whose value determines which hyperbola slice
        is displayed.  Typical examples: stress ratio, temperature, frequency.
    \item[Vertical asymptote]
        The value $X_\text{asym}$ at which the fitted hyperbola becomes undefined
        (denominator $\to 0$).
\end{description}

% ══════════════════════════════════════════════════════════════════════════════

\vfill
\begin{center}
\color{gray}\small
\rule{0.6\linewidth}{0.4pt}\\[4pt]
Fatigue 3D Model --- User Manual\\
Document prepared automatically from source code analysis.
\end{center}

\end{document}
